\section{API végpontok és adatáramlások}

A WatchDB a Next.js API route-jait használja a kliensoldali műveletek kiszolgálására. A böngésző nem közvetlenül kommunikál sem a Supabase adatbázissal, sem a TMDB API-val: ezek elérése szerveroldalon történik. Ez egyrészt védi a külső szolgáltatásokhoz tartozó kulcsokat, másrészt egységesebb hibakezelést és válaszformátumot tesz lehetővé. A sikeres válaszok \texttt{data} mezőben, a hibák pedig \texttt{error} mezőben térnek vissza.

\begin{table}[H]
\centering
\scriptsize
\begin{tabular}{|p{0.22\textwidth}|p{0.14\textwidth}|p{0.27\textwidth}|p{0.27\textwidth}|}
\hline
\textbf{Végpont} & \textbf{Metódus} & \textbf{Paraméterek} & \textbf{Szerep} \\
\hline
\texttt{/api/search} & \texttt{GET} & \begin{tabular}[t]{@{}l@{}}\texttt{query}: kötelező\\\texttt{page}: opcionális\end{tabular} & Filmkeresés a TMDB keresési endpoint-ján keresztül. \\
\hline
\texttt{/api/watchlist} & \texttt{GET} & \begin{tabular}[t]{@{}l@{}}Nincs query.\\Felhasználó: sessionből.\end{tabular} & A bejelentkezett felhasználó watchlistjének lekérése. \\
\hline
\begin{tabular}[t]{@{}l@{}}\texttt{/api/watchlist/}\\\texttt{[movieId]}\end{tabular} & \begin{tabular}[t]{@{}l@{}}\texttt{GET}\\\texttt{POST}\\\texttt{DELETE}\end{tabular} & \begin{tabular}[t]{@{}l@{}}\texttt{movieId}: route paraméter\\Nincs request body.\end{tabular} & Egy film watchlist állapotának lekérése, hozzáadása vagy eltávolítása. \\
\hline
\texttt{/api/watched} & \texttt{GET} & \begin{tabular}[t]{@{}l@{}}\texttt{page}: opcionális\\\texttt{limit}: opcionális\end{tabular} & A megnézett filmek paginált listájának lekérése. \\
\hline
\begin{tabular}[t]{@{}l@{}}\texttt{/api/watched/}\\\texttt{[movieId]}\end{tabular} & \begin{tabular}[t]{@{}l@{}}\texttt{GET}\\\texttt{POST}\\\texttt{DELETE}\end{tabular} & \begin{tabular}[t]{@{}l@{}}\texttt{movieId}: route paraméter\\POST body: \texttt{rating}\\POST body: \texttt{watchedAt}\end{tabular} & Egy film watched állapotának lekérése, megnézettként jelölése vagy eltávolítása. \\
\hline
\end{tabular}
\caption{A WatchDB fő API végpontjai}
\label{tab:watchdb-api-endpoints}
\end{table}

A lekérdező végpontok kétféle adatforrással dolgoznak. A keresés és a filmrészletek a TMDB API-ból származó adatokat használják, míg a watchlist és watched listák először a Supabase adatbázisból kérik le a felhasználóhoz tartozó sorokat. A visszakapott \texttt{tmdb\_movie\_id} értékek alapján az alkalmazás ezután lekéri a szükséges filmadatokat a TMDB-ből, majd a saját adatbázisból és a külső API-ból származó adatokat egy közös válaszobjektumban adja vissza.

Az állapotmódosító műveleteknél a kliensoldalon a TanStack Query kezeli a lekérdezések és mutációk (\texttt{CREATE, UPDATE, DELETE}) állapotát. Egy watchlist vagy watched művelet után az érintett query-k frissülnek, így a filmrészletek oldalon és a profiloldalon megjelenő adatok újratöltés nélkül is követhetik a változást.

\begin{figure}[H]
\centering
\scriptsize
\resizebox{\textwidth}{!}{%
\begin{tikzpicture}[
    actor/.style={
        draw,
        rounded corners,
        fill=gray!8,
        align=center,
        minimum height=0.68cm,
        text width=3cm
    },
    message/.style={-{Latex[length=2mm]}, thick},
    returnmessage/.style={-{Latex[length=2mm]}, thick, dashed},
    lifeline/.style={gray, dashed},
    msglabel/.style={
        fill=white,
        inner sep=1.5pt,
        align=center
    },
    note/.style={
        draw,
        rounded corners,
        fill=gray!8,
        align=center,
        text width=3.2cm,
        inner sep=4pt
    }
]
\node[actor] (client) at (0,0.5) {Kliens\\filmoldal};
\node[actor] (api) at (4.2,0.5) {Next.js\\API route};
\node[actor] (auth) at (8.4,0.5) {Supabase\\Auth};
\node[actor] (database) at (12.6,0.5) {Supabase\\adatbázis};

\draw[lifeline] (0,0.15) -- (0,-10.8);
\draw[lifeline] (4.2,0.15) -- (4.2,-10.8);
\draw[lifeline] (8.4,0.15) -- (8.4,-10.8);
\draw[lifeline] (12.6,0.15) -- (12.6,-10.8);

\draw[fill=white] (-0.12,-0.75) rectangle (0.12,-9.75);
\draw[fill=white] (4.08,-1.1) rectangle (4.32,-9.05);
\draw[fill=white] (8.28,-3.75) rectangle (8.52,-4.55);
\draw[fill=white] (12.48,-5.65) rectangle (12.72,-8.25);

\draw[message] (0.12,-1.25) -- node[msglabel, above] {1. \texttt{POST /api/watched/[movieId]}} node[msglabel, below] {body: \texttt{rating}, \texttt{watchedAt}} (4.08,-1.25);
\draw[message] (4.32,-2.15) -- (5.15,-2.15) -- (5.15,-2.85) -- node[msglabel, right] {2. bemenet\\validálása} (4.32,-2.85);
\draw[message] (4.32,-4.0) -- node[msglabel, above] {3. \texttt{auth.getUser()}} (8.28,-4.0);
\draw[returnmessage] (8.28,-4.55) -- node[msglabel, above] {4. \texttt{user\_id} vagy hiba} (4.32,-4.55);
\draw[message] (4.32,-5.7) -- node[msglabel, above] {5. \texttt{watched.insert(...)}} node[msglabel, below] {\texttt{rating}, \texttt{watched\_at}} (12.48,-5.7);
\draw[returnmessage] (12.48,-6.55) -- node[msglabel, above] {6. mentett watched sor} (4.32,-6.55);
\draw[message] (4.32,-7.5) -- node[msglabel, above] {7. \texttt{watchlist.delete(...)}} node[msglabel, below] {ha szerepelt a listán} (12.48,-7.5);
\draw[returnmessage] (12.48,-8.25) -- node[msglabel, above] {8. törlés eredménye} (4.32,-8.25);
\draw[returnmessage] (4.08,-9.15) -- node[msglabel, above] {9. \texttt{201 Created}} node[msglabel, below] {friss watched állapot} (0.12,-9.15);
\node[note] at (0,-10.3) {10. \texttt{watched} és \texttt{watchlist} query-k frissítése};
\end{tikzpicture}%
}
\caption{A megnézettként jelölés adatáramlása}
\label{fig:watchdb-watched-flow}
\end{figure}
